\documentclass[tg]{mdtufsm}

\usepackage{caption}
\captionsetup{
	justification=raggedright,
	figurewithout,labelsep=endash,
	singlelinecheck=false
	}





\usepackage{chngcntr}
\counterwithout{figure}{chapter}
\counterwithout{table}{chapter}
\usepackage{ragged2e}
\usepackage{colortbl}
\usepackage{times}              % pacote para usar fonte Adobe Times
\usepackage{setspace}
\usepackage{array,multirow}
\usepackage{float} 
\usepackage[portuguese,titlenumbered,ruled]{algorithm2e} % algoritmos
\usepackage[T1]{fontenc}        % pacote para conj. de caracteres correto
\usepackage{fix-cm} %para funcionar corretamente o tamanho das fontes da capa
\usepackage{times, color, xcolor}       % pacote para usar fonte Adobe Times e cores
\usepackage[utf8]{inputenc}   % pacote para acentuação
\usepackage{graphicx}  % pacote para importar figuras
\usepackage{amsmath,latexsym,amssymb} %Pacotes matemáticos
\usepackage[%hidelinks%, 
            bookmarksopen=true,linktocpage,colorlinks=true,
            linkcolor=black,citecolor=black,filecolor=magenta,urlcolor=blue,
            pdftitle={Título do Trabalho},
            pdfauthor={nome do autor},
            pdfsubject={Trabalho de Conclusão de Curso},
            pdfkeywords={modelo, latex, tcc, graduação}
            ]{hyperref} 

%Margens conforme MDT 2015
\usepackage[inner=30mm,outer=20mm,top=30mm,bottom=20mm]{geometry} 

%==============================================================================
% Se o pacote hyperref foi carregado a linha abaixo corrige um bug na hora
% de montar o sumário da lista de figuras e tabelas
% Se o pacote não foi carregado, comentar a linha %
%==============================================================================
\input{macros/bugcaption}

%==============================================================================
% Identificação do trabalho
%==============================================================================
\title{Template para TCC IFFar - SVS}

\author{Sobre nome}{nome do autor}
%Descomentar se for uma "autora"
%\autoratrue

%% Lombada
\titulacao{Tecnólogo em Análise e Desenvolvimento de Sistemas} %Tecn\'{o}logo ou Bacharel ou Licenciatura



%folha rosto
\minedu{MINISTÉRIO DA EDUCAÇÃO}
\cabecalho{SECRETARIA DE EDUCAÇÃO PROFISSIONAL E TECNOLÓGICA INSTITUTO FEDERAL DE EDUCAÇÃO, CIÊNCIA E TECNOLOGIA FARROUPILHA - CAMPUS SÃO VICENTE DO SUL}
\course{Curso Superior de Tecnologia em Análise e Desenvolvimento de Sistemas}
\initials{ADS} 
\degree{Tecnólogo em Análise e Desenvolvimento de Sistemas}

\faculdade{Instituto Federal Farroupilha - Campus São Vicente do Sul}
\faculdadeing{Campus São Vicente do Sul - Instituto Federal Farroupilha}
\usepackage{hyphenat} % evitar separação
\siglafaculdade{\nohyphens{IFFar - SVS}, RS}
\campus{São Vicente do Sul}

\cidade{São Vicente do Sul}
\siglaestado{RS}
\pais{Brasil}


% Número do TG (verificar na secretaria do curso)
% Para mestrado deixar sem opção dentro do {}
\trabalhoNumero{}

%Orientador
\advisoronepage[Prof.]{Me.}{Rissetti}{Gustavo}
\advisor[Prof.]{Me. (IFFar - SVS)}{Rissetti}{Gustavo}
%Se for uma ``orientadora'' descomentar a linha baixo
%\orientadoratrue

%Co orientador, comentar se não existir
\coadvisoronepage[Prof.]{Me.}{Rissetti}{Gustavo}
\coadvisor[Prof.]{Me. (IFFar - SVS)}{Rissetti}{Gustavo}
\coorientacaotrue %Se houver coorientação
%\coorientadoratrue %Se for uma ``Co-Orientadora''

%Avaliadores (Banca)
\committee[Me.]{Sobre nome}{Nome menbro banca}{UFSM}
\committee[Tecg.]{Sobre nome}{Nome menbro banca}{UFSM}

% a data deve ser a da defesa; 
% dia, mes e ano correntes
\date{28}{03}{2018} 

%Palavras chave
\keyword{modelo}
\keyword{latex} 
\keyword{tcc}
\keyword{graduação}

%keyword
\keywordAbstract{model}
\keywordAbstract{latex}
\keywordAbstract{tcc}
\keywordAbstract{graduation}
%%=============================================================================
%% Início do documento
%%=============================================================================
\begin{document}
	
	

%%=============================================================================
%% Capa e folha de rosto
%%=============================================================================
\maketitle
%%=============================================================================
%% Catalogação (obrigatório para mestrado) e Folha de aprovação
%%=============================================================================
%Somente obrigatório para dissertação, para TG, remover as linhas	77	%
%Como a CIP vai ser impressa atrás da página de rosto, as margens inner e outer	
%devem ser invertidas.

\newgeometry{inner=20mm,outer=30mm,top=30mm,bottom=20mm}	
\makeCIP{autor@email} %email do autor		
\restoregeometry

%Se for usar a catalogação gerada pelo gerador do site da biblioteca comentar as linhas
%\includeCIP{CIP.pdf}
	
%%=============================================================================
%folha de aprovação
%%=============================================================================
\makeapprove

%%=============================================================================
%% Dedicatória (opcional)
%%=============================================================================
\include{pessoal/dedicatoria}
%%=============================================================================
%% Agradecimentos (opcional)
%%=============================================================================
\include{pessoal/agradecimento}
%%=============================================================================
%% Epígrafe (opcional)
%%=============================================================================
\include{pessoal/epigrafe}
%%=============================================================================
%% Resumo
%%=============================================================================
\begin{abstract}
		
Aqui você escreve o resumo. Lembrando no máximo 250 palavras para tcc e 500 palavras para tese ou dissertação.

\end{abstract}

%%=============================================================================
%% Abstract
%%=============================================================================
% resumo na outra língua
% como parametros devem ser passados o titulo, o nome do curso,
% as palavras-chave na outra língua, separadas por vírgulas, o mês em inglês
%o a sigla do dia em inglês: st, nd, th ...
\begin{englishabstract}
{Abstract title}
%%============ MANTER ESSA LINHA ============================================
\ %%============ MANTER ESSA LINHA ==========================================
%%============ MANTER ESSA LINHA ============================================
\ %%============ MANTER ESSA LINHA ==========================================
%%============ MANTER ESSA LINHA ============================================
\ %%============ MANTER ESSA LINHA e deixar uma em branco ===================


Here you write the summary. Remembering a maximum of 250 words for tcc and 500 words for thesis or dissertation.

\end{englishabstract}

%% Lista de Ilustrações (opc)
%% Lista de Símbolos (opc)
%% Lista de Anexos e Apêndices (opc)

%%=============================================================================
%% Lista de figuras (comentar se não houver)
%%=============================================================================
\listoffigures
%%=============================================================================
%% Lista de tabelas (comentar se não houver)
%%=============================================================================
\listoftables

%%=============================================================================
%% Lista de Apêndices (comentar se não houver)
%%=============================================================================
%\listofappendix

%%=============================================================================
%% Lista de Anexos (comentar se não houver)
%%=============================================================================
%\listofannex

%%=============================================================================
%% Lista de abreviaturas e siglas
%%=============================================================================
 %o parametro deve ser a abreviatura mais longa

\include{siglas/lista}

%%=============================================================================
%% Lista de simbolos (opcional)
%%=============================================================================
%Simbolos devem aparecer conforme a ordem em que aparecem no texto
% o parametro deve ser o símbolo mais longo
%\begin{listofsymbols}{teste}
%  \item [$\varnothing$] vazio
 % \item [$\Gamma$]  Gama
%  \item [$\forall$] Para todo
%\end{listofsymbols}

%%=============================================================================
%% Sumário
%%=============================================================================
\tableofcontents











%%=============================================================================
%% Início da dissertação
%%=============================================================================
\setlength{\baselineskip}{1.5\baselineskip}

%Adiciona cada capitulo
\include{capitulos/introducao}
\include{capitulos/desenvolvimento}
\include{capitulos/conclusao}

\setlength{\baselineskip}{\baselineskip}

%%=============================================================================
%% Referências
%%=============================================================================
\begin{flushleft}
\bibliographystyle{abnt}
\bibliography{referencias/referencias}
\end{flushleft}

%IMPORTANTE: Se precisar usar alguma seção ou subseção dentro dos apêndices ou
%anexos, utilizar o comando \tocless para não adicionar no Sumário
%Exemplos: 
% \tocless\section{Histórico}
%%=============================================================================
%% Apêndices
%%=============================================================================
\appendix

\include{apendices/apendicea}
%\include{apendices/apendiceb}

%%=============================================================================
%% Anexos
%%=============================================================================
\annex
\include{anexos/anexoa}

\end{document}
